\author{OTS}
\documentclass[14pt]{../../inc/mybib}

\setincpath{../../inc/}

\usepackage{bible_style}
\usepackage{textmakro}
\graphicspath{{../../assets/images/}}
\usepackage{header}

\newcommand{\Name}{Fredy}

% ensure scrlayer-scrpage has sufficient footheight
\setlength{\footheight}{10.4pt}
\setlength{\headheight}{0.4pt}

\begin{document}
\begin{bibelbox}{SCHL}{Kol}{3:15-17}
    Lasst das Wort des Christus reichlich in euch wohnen in aller Weisheit; lehrt und ermahnt einander und singt mit Psalmen und Lobgesängen und geistlichen Liedern dem Herrn lieblich in eurem Herzen. Und was immer ihr tut in Wort und Werk, das tut alles im Namen des Herrn Jesus und dankt Gott, dem Vater, durch ihn.
\end{bibelbox}
\section{Begrüssung}
Ich möchte auch euch alle recht herzlich zum Gottesdienst begrüssen. Schön, dass ihr gekommen seid, um unserem \herr N zu danken, ihn zu loben und zu preisen.

Auch dich \Name{} will ich herzlich begrüssen. Schön, dass du wieder bei uns bist und mit uns Gottesdienst feierst.
\begin{itemize}
    \item[] \beten{}
    \item[] Und anschliessend singen wir zusammen das Lied
    \item[] \lied{Wunderbar. grosser Erlöser}.
\end{itemize}

\section{Informationen}
\begin{itemize}
    \item \bt[infos]{Bibel und Gebetsabend:} Do, 19.03.2026 20:00 Uhr Taufe in Dübendorf mit Fredy Peter. Wir werden hier die Zeugnisse übertragen und die Möglichkeit bieten für die Täuflinge und andere Anliegen zu beten.
    \item \bt[infos]{Nächster Gottesdienst:} So, 22.03.2026 14:45 Uhr mit Nathanael Winkler
    \item \bt[infos]{Ostern:} Ostern von 5. April findet hier kein Gottesdienst statt. Ihr seid alle herrzlich eingeladen an der Osterkonferenz in Dübendorf teil zu nehmen.
    \item \bt[infos]{Vortrag:} Am 13. April wird uns Nathanel Winkler im Rahmen seiner Schweizer Turnee ihr in diesen Räumen besuchen und einen Vortrag über Gog und Magog: Alte Prophezeiung, neue Geopolitik halten. Ladet Freunde und Kollegen ein.
    \item \bt[infos]{Kollekte:} Die Kollekte wird hinten in der grünen Box gesammelt und wird vollumfänglich für den Bau dieser Gemeinde eingesetzt.
\end{itemize}

\section{ Input }
\begin{spacing}{1.5}
    \begin{block}[Das Tausendjährige Reich]
        Ich will euch aber trotzdem nochmals die wichtigsten Bibelverse vorlesen.
        \bibleverse{Offb}(20:1-10) 
        Thomas hat uns hier zwei Sonntage das Tausendjährige Reich nahe gebracht. Darum brauche ich euch heute nicht damit zu langweilen sondern dies als Anlass nehmen etwas über Irrlehren zu sagen. Wie auch schon Thomas gesagt hat, glaube auch ich, dass diese 1000 Jahre wörtlich zu nehmen sind. Wieso sollte diese Zahl sonst so oft erwähnt werden? Manche Ausleger sagen, dass die Zahl 1000 Vollkommenheit bedeutet. Diese Sichtweise wird Amillennialismus genannt, also \enquote{kein Millenium}. Eine Theorie ist auch, dass wir jetzt schon im Millenium sind. Die kath. Kirche vertritt diese Ansicht. Sie sagt, Jesus kommt wieder, wenn hier auf der Erde Friede und Ruhe herrscht.

        Andere Theorien ist, dass das tausendjährige Reich mit dem ersten kommen Jesus begonnen hat. Das Königreich ist schon errichtet und jeder der Christus von Herzen aufnimmt, trägt das Königreich in sich. Der Teufel wurde mit dem Tod und der Auferstehung Christi gebunden. Es sind ziemlich bekannte YouTube Prediger, die diese Ansicht vertreten. So zB auch Olaf Latzel, Tobias Riemenscheider, Willy Zorn, Jürgen Fischer.
        Ich habe auf Youtube eine Debatte gehört von einem Millenialist und einem Amillienalist. Beides Bibeltreue Christen und beide haben mit der Bibel argunmentiert, aber als Amillienalist muss man dann halt schon viele Bibelstellen vergeistlichen. Für mich ist darum, an ein buchstäbliches Millenium zu glauben, die bei weitem einfachere Methode. Wir sollten aber immer beachten, dass nur weil jemand zu diesem Thema eine andere Sicht hat, ist er darum noch kein Irrlehrer. Irrlehrer sind solche die das falsche Evangelium verkünden.   
        \textbf{Biblische Merkmale für Irrlehrer}
        \begin{enumerate}
            \item \textbf{Falsches Evangelium}: Wer ein anderes Evangelium bringt als das apostolische, steht unter scharfer Warnung (Galater 1,6-9).
            \item \textbf{Falsches Falsches Christusbild}: Wer Jesus nicht gemäß der biblischen Wahrheit bekennt (Menschwerdung, Herrschaft, Werk), gilt als „nicht aus Gott“ (1. Johannes 4,1-3, 2. Johannes 7).
            \item \textbf{Verdrehung der Gnade:} Gnade wird als Freibrief zur Sünde missbraucht (Judas 4).
            \item \textbf{Leugnung zentraler Glaubensinhalte:} z. B. Auferstehung oder Herrschaft Christi (2. Timotheus 2,17-18, 2. Petrus 2,1).
            \item \textbf{Schlechter Fruchtcharakter:} Habgier, Ausbeutung, Stolz, Unmoral, Spaltung (Matthäus 7,15-20, 2. Petrus 2,2-3, Titus 3,10).
        \end{enumerate}
        Bevor Jesus in Macht und Herrlichkeit wieder kommt, werden viele solche Irrlehrer auftreten. In der heutigen Zeit mit dieser weltweiten Vernetzung ist es einfach möglich ein falsches Evangelium zu verbreiten und es ist leider nicht immer so offensichtlich. Vieles wird in die Liebe eingepackt, in die Gnade und an den barmherzigen Gott der doch alle liebt.

        Ja Gott liebt alle Menschen, er liebt dich und mich. Gott hat aber auch Gefühle, er wird zornig, kann entäuscht werden, ist traurig und das nur wegen uns Menschen die ihn dauernd enttäuschen. Gott liebt uns so sehr, dass sein Sohn für uns am Kreuz gestorben ist. Wie Danken wir ihm dafür? Er verlangt ja von nur, dass wir ihn lieben mehr nicht. Lieben wir Gott? Oder machen wir das alles nur um in den Himmel zu kommen?

        Und für alle die denken, dass der Mensch doch im Grunde seines Herzens gut ist, sehen wir hier im Millenium der Beweis, dass es nicht so ist. Satan ist Gebunden im Abgrund. 
        \begin{block}[Zitat von Warren W. Wiersbe:]
            In gewissem Sinne wird das tausendjährige Reich all das \enquote{zusammenfassen}, was Gott über das Herz des Menschen in den verschiedenen Geschichtsabschnitten gesagt hat. Es wird eine Herrschaft des Gesetzes sein, und doch wird das Gesetz das sündige Herz des Menschen nicht verändern können. Der Mensch wird sich immer noch gegen Gott auflehnen.
        \end{block}        
        Er kann nichts machen und obwohl keine Krankheiten, kein Hunger, keine Kriege mehr gibt wenden sich Menschen von Jesus ab. Auch perfekte Bedingungen bringen keine guten Herzen hervor. Darum wird Gott zum Abschluss den Teufel nochmals los lassen. Der Satan wird diese Abtrünnigen um sich scharen um gegen Jesus anzutreten. Gott wird aber dieses mal kurzen Prozess machen Off 20:7-10. Das Böse wird vernichtet. Jetzt wirds spannend. Was passiert wenn das Böse vernichtet wird? Das hört ihr dann am nächsten Sonntag.   
        
    \end{block}
    \lied{Wir wissen nicht den Tag noch die Stunde}.
\end{spacing}
\section{Predigt}
\textbf{Nach der Predigt}\newline
Danken für die Predigt.
% \section{Abendmahl}
% \begin{itemize}
%     \item [] Beten für das Brot
%     \item [] Beten für den Wein
% \end{itemize}
\section{Abschluss}
\begin{itemize}
    \item [] \lied{O Gott dir sei Ehre}.
    \item [] \beten{}
\end{itemize}
\begin{bibelbox}{SCHL}{IIKor}{13:13}
    Die Gnade des Herrn Jesus Christus und die Liebe Gottes und die Gemeinschaft des Heiligen Geistes sei mit euch allen! Amen
\end{bibelbox}
\end{document}